\documentclass{../yalumba}

\title{Pair-Coupled Ecosystem Transport for\\Reference-Free Relatedness Detection:\\Spectral Ecology v3}
\author{Ajax Davis}
\date{April 2026}

\setabstract{%
We introduce \emph{pair-coupled ecosystem transport}, the first cross-sample
structural comparison method in the yalumba symbiogenesis program.  Where
Spectral Ecology v2 compared genomes through separate per-sample spectral
summaries---hitting a ceiling at $+1.55\%$ separation because within-sample
invariants cannot distinguish ``similar because related'' from ``similar
because similar coverage''---v3 builds a structure that \emph{only exists when
two samples are considered together}: a bipartite soft correspondence between
their module ecosystems.  For each pair of samples, we (1)~construct enriched
20-dimensional node feature vectors per module combining intrinsic properties
(support, cohesion, motif count, clustering coefficient), spectral properties
(heat kernel signatures at 5 timescales, spectral position from 4
eigenvectors), and ecological role (5-dimensional one-hot), all
rank-normalized within each sample; (2)~compute a cross-sample Gaussian
similarity kernel $K(m_A, m_B) = \exp(-d(x_A, x_B)/\tau)$ with auto-scaled
bandwidth; (3)~extract mutual top-5 bipartite matches (module $i$ in $A$
lists $j$ in $B$ among its top~5 AND $j$ lists $i$ among its top~5); and
(4)~score via three components: transport affinity (total normalized
compatibility mass, weight 0.45), topology preservation (fraction of matched
neighbors that map to each other's neighbors, weight 0.40), and unmatched
penalty ($1 - \text{fraction unmatched}$, weight 0.15).  On the CEPH~1463
six-member pedigree benchmark, v3 achieves $+4.93\%$ mean separation---a
$3.2\times$ improvement over v2's $+1.55\%$ and competitive with Module
Persistence v3's $+3.36\%$.  All four parent--child pairs rank above all
four pure unrelated pairs for the first time in any spectral method.  Two
parent--child pairs (Mat.GM~$\to$~Mother at 1.108, Pat.GM~$\to$~Father at
1.103) rank 2nd and 3rd.  The weakest gap remains negative ($-54.67\%$)
because the spouse pair (Father~$\leftrightarrow$~Mother) scores 1.477---
the transport affinity component is not yet selective enough, finding
1{,}400 mutual matches between graphs of 444 and 443 modules.  We diagnose
this as an indiscriminate matching problem: when nearly every module finds
a plausible mate, the transport score reflects graph size rather than
structural specificity.  The fix is clear: penalize matches that are not
significantly better than random, or require matched neighborhoods (not just
matched nodes) to be coherent.  The architectural shift from within-sample
summaries to cross-sample correspondence is validated: the $3.2\times$
separation improvement demonstrates that pair-specific structure is the
missing ingredient the symbiogenesis program has been seeking.}

\begin{document}
\maketitle

\tableofcontents
\newpage


% ════════════════════════════════════════════════════════════════════
\section{Introduction}
\label{sec:intro}

\subsection{The Within-Sample Ceiling}

Spectral Ecology v2 demonstrated that identity-free spectral methods can
detect real organizational similarity between related individuals, achieving
$+1.55\%$ separation on the CEPH~1463 pedigree.  Five iterative experiments
(co-occurrence graph construction, eigenvalue filtering, node-level HKS
distribution comparison, mean normalization) converged to a ceiling of
$\sim$1.5--1.9\% regardless of which invariants were used.

\ymarginnote{v2 proved the signal exists.  v3 extracts it.}

The ceiling exists because all v2 features are \emph{within-sample summaries}:
eigenvalue density, heat kernel trace, role spectrum, localization statistics.
Two samples can have identical summaries for two reasons: (1)~they share
inherited genomic organization, or (2)~they have similar sequencing depth and
quality.  Within-sample methods cannot distinguish these cases.

The v2 paper concluded: ``further progress requires cross-sample structural
signal.''  This paper delivers it.

\subsection{The Cross-Sample Insight}

The key insight is simple: instead of separately summarizing each sample and
comparing the summaries, build a structure that \emph{only exists when two
samples are considered together}.

\begin{pullquote}
Can the modules in sample $A$ be cheaply explained as transformed versions
of the modules in sample $B$?
\end{pullquote}

If two individuals are related, their module ecosystems should admit many
coherent local correspondences: module $i$ in $A$ matches module $j$ in $B$
(similar ecological role, similar diffusion profile, similar neighborhood
structure), and $i$'s neighbors also match $j$'s neighbors.  If two
individuals are unrelated, correspondences should be sparse and incoherent.

This is fundamentally a \emph{transport} question: how much structural mass
can be moved from one ecosystem to the other at low cost?

\subsection{Contributions}

\begin{enumerate}
  \item \textbf{20-dimensional enriched node features} combining intrinsic,
    spectral, and ecological properties, rank-normalized within each sample.
  \item \textbf{Mutual top-$k$ bipartite matching}: a pair-specific structure
    that captures cross-sample module correspondence without requiring exact
    identity.
  \item \textbf{Three-component scoring}: transport affinity (total mass),
    topology preservation (neighborhood coherence), and unmatched penalty
    (coverage of the ecosystem).
  \item \textbf{$3.2\times$ separation improvement}: $+4.93\%$ vs.\
    $+1.55\%$, validating cross-sample structure as the missing ingredient.
  \item \textbf{All parent--child pairs above all pure unrelated pairs}
    in the ranking---a first for spectral methods.
\end{enumerate}


% ════════════════════════════════════════════════════════════════════
\section{Methods}
\label{sec:methods}

\subsection{Per-Sample Preparation (Unchanged from v2)}

For each sample independently:

\begin{enumerate}
  \item \textbf{Module extraction.}  Raw reads $\to$ k-mer co-occurrence
    graphs $\to$ connected components.  Parameters: $k=15$, window $=150$\,bp,
    min support $=3$, min cohesion $=0.25$.
  \item \textbf{Co-occurrence graph.}  Read-level module co-occurrence with
    normalized weights:
    $w(A,B) = \text{cooccur}(A,B) / \sqrt{\text{freq}(A) \cdot \text{freq}(B)}$.
  \item \textbf{Spectral decomposition.}  Normalized Laplacian $\to$ Jacobi
    eigendecomposition $\to$ heat kernel at 10 timescales $\to$ role assignment.
\end{enumerate}

This produces 295--449 modules per sample with 611--1{,}410 co-occurrence
edges.

\subsection{Enriched Node Features}

Each module receives a 20-dimensional feature vector:

\begin{table}[H]
  \centering
  \tablestyle
  \caption{Node feature vector (20 dimensions).  All features are
    rank-normalized within each sample to remove absolute scale.}
  \label{tab:features}
  \begin{tabular}{clc}
    \toprule
    \rowcolor{table-header}
    \textcolor{white}{\textbf{Dims}} &
    \textcolor{white}{\textbf{Feature}} &
    \textcolor{white}{\textbf{Source}} \\
    \midrule
    1 & $\log(\text{support} + 1)$ & Module extraction \\
    1 & Cohesion & Module extraction \\
    1 & $\log(\text{motif count} + 1)$ & Module extraction \\
    1 & Weighted degree & Co-occurrence graph \\
    1 & Clustering coefficient & Co-occurrence graph \\
    5 & HKS at $t \in \{0.1, 0.5, 2, 10, 50\}$ & Heat kernel \\
    4 & Spectral position ($\phi_1 \ldots \phi_4$) & Eigenvectors \\
    5 & Role one-hot (C/S/B/F/Y) & Role assignment \\
    \midrule
    \textbf{20} & \textbf{Total} & \\
    \bottomrule
  \end{tabular}
\end{table}

\ymarginnote{Rank normalization removes coverage confound at the feature level.}

\begin{definition}
\textbf{Rank normalization.}  For each feature dimension, values are replaced
by their within-sample rank divided by $n-1$, producing values in $[0, 1]$.
Tied values receive the average rank.  This is robust to outliers and removes
absolute scale---a module with degree 50 in a sparse graph and a module with
degree 200 in a dense graph both receive rank $\sim$0.95 if they are
top-5\% in their respective ecosystems.
\end{definition}

\subsection{Cross-Sample Similarity Kernel}

For samples $A$ and $B$ with module feature matrices $X_A \in \mathbb{R}^{n_A
\times 20}$ and $X_B \in \mathbb{R}^{n_B \times 20}$, define the pairwise
similarity kernel:

\begin{equation}
  K(i, j) = \exp\!\left(-\frac{\|x_i^A - x_j^B\|_2}{\tau}\right)
  \label{eq:kernel}
\end{equation}

where $\tau$ is the median of all pairwise Euclidean distances (auto-scaled
bandwidth).  This produces an $n_A \times n_B$ matrix of soft
compatibilities.

\subsection{Mutual Top-$k$ Bipartite Matching}

\begin{definition}
A \textbf{mutual match} $(i, j)$ exists when module $i$ in $A$ lists module
$j$ in $B$ among its top-$k$ most similar modules AND module $j$ in $B$
lists module $i$ in $A$ among its top-$k$.  We use $k = 5$.
\end{definition}

\ymarginnote{Mutual matching is more selective than one-directional.}

Mutual matching is strictly more selective than one-directional: a module must
be among the best candidates from \emph{both} sides.  This suppresses
spurious matches where a low-quality module in one sample happens to be
near a high-quality module in the other.

\subsection{Three-Component Scoring}

\subsubsection{Component 1: Transport Affinity (weight 0.45)}

The total normalized compatibility mass across all mutual matches:

\begin{equation}
  T(A, B) = \frac{\sum_{(i,j) \in \mathcal{M}} K(i, j)}
                  {\sqrt{n_A \cdot n_B}}
  \label{eq:transport}
\end{equation}

where $\mathcal{M}$ is the set of mutual matches.  Higher $T$ means more
modules find high-quality mates in the other ecosystem.

\subsubsection{Component 2: Topology Preservation (weight 0.40)}

For each mutual match $(i, j)$, measure whether their neighborhoods are
also matched:

\begin{equation}
  N(i, j) = 0.6 \cdot \frac{|\{\text{matched neighbors of } i \text{ that map to neighbors of } j\}|}
                              {|\text{matched neighbors of } i|}
           + 0.4 \cdot \frac{\min(\deg(i), \deg(j)) + 1}
                             {\max(\deg(i), \deg(j)) + 1}
  \label{eq:topology}
\end{equation}

The topology preservation score is the average $N(i,j)$ over all mutual
matches.  High topology preservation means the matching respects graph
structure, not just node features.

\ymarginnote{Topology preservation is the key: it tests whether the matching
is structurally coherent, not just feature-similar.}

\subsubsection{Component 3: Unmatched Penalty (weight 0.15)}

\begin{equation}
  U(A, B) = \frac{|\text{matched}_A| + |\text{matched}_B|}{n_A + n_B}
  \label{eq:unmatched}
\end{equation}

The fraction of modules that found a mutual mate.  If many modules are
unmatched, the ecosystems are less compatible.

\subsubsection{Composite}

\begin{equation}
  S(A, B) = 0.45 \cdot T + 0.40 \cdot N + 0.15 \cdot U
  \label{eq:composite}
\end{equation}


% ════════════════════════════════════════════════════════════════════
\section{Results}
\label{sec:results}

\subsection{Full Score Table}

\begin{table}[H]
  \centering
  \tablestyle
  \caption{Complete pairwise scores for Spectral Ecology v3, ranked by
    composite similarity.  Parent--child pairs marked with
    $\blacktriangleright$.  $T$ = transport affinity, $N$ = topology
    preservation, $U$ = unmatched penalty, $|\mathcal{M}|$ = mutual match
    count.}
  \label{tab:scores}
  \begin{tabular}{clccccl}
    \toprule
    \rowcolor{table-header}
    \textcolor{white}{\textbf{Rank}} &
    \textcolor{white}{\textbf{Pair}} &
    \textcolor{white}{\textbf{$S$}} &
    \textcolor{white}{\textbf{$T$}} &
    \textcolor{white}{\textbf{$N$}} &
    \textcolor{white}{\textbf{$|\mathcal{M}|$}} &
    \textcolor{white}{\textbf{Type}} \\
    \midrule
    1 & Father $\leftrightarrow$ Mother
      & 1.477 & 2.640 & 0.375 & 1{,}400 & Spouses \\
    2 & $\blacktriangleright$ Mat.GM $\to$ Mother
      & 1.108 & 1.910 & 0.367 & 855 & Parent--child \\
    3 & $\blacktriangleright$ Pat.GM $\to$ Father
      & 1.103 & 1.837 & 0.388 & 829 & Parent--child \\
    4 & Mat.GF $\leftrightarrow$ Father
      & 1.011 & 1.628 & 0.372 & 843 & In-law \\
    5 & Pat.GF $\leftrightarrow$ Mother
      & 0.954 & 1.573 & 0.337 & 651 & In-law \\
    6 & $\blacktriangleright$ Pat.GF $\to$ Father
      & 0.943 & 1.487 & 0.382 & 636 & Parent--child \\
    7 & $\blacktriangleright$ Mat.GF $\to$ Mother
      & 0.930 & 1.380 & 0.411 & 820 & Parent--child \\
    8 & Pat.GM $\leftrightarrow$ Mat.GM
      & 0.903 & 1.470 & 0.335 & 704 & Unrelated \\
    9 & Pat.GF $\leftrightarrow$ Mat.GM
      & 0.746 & 1.112 & 0.332 & 471 & Unrelated \\
    10 & Pat.GF $\leftrightarrow$ Mat.GF
      & 0.741 & 1.138 & 0.359 & 677 & Unrelated \\
    \bottomrule
  \end{tabular}
\end{table}

\subsection{Separation Analysis}

\begin{table}[H]
  \centering
  \tablestyle
  \caption{Summary metrics for Spectral Ecology v3.}
  \label{tab:summary}
  \begin{tabular}{lr}
    \toprule
    \rowcolor{table-header}
    \textcolor{white}{\textbf{Metric}} &
    \textcolor{white}{\textbf{Value}} \\
    \midrule
    Parent--child average & 1.021 \\
    Unrelated average & 0.972 \\
    Mean separation & $+4.93\%$ \\
    Weakest gap & $-54.67\%$ \\
    Detects all & NO \\
    Prep time & 741\,s \\
    Compare time & 1.5\,s (0.15\,s/pair) \\
    Heap delta & $+$1{,}150\,MB \\
    \bottomrule
  \end{tabular}
\end{table}

\begin{keyfinding}
\textbf{Separation tripled}: $+4.93\%$ (v3) vs.\ $+1.55\%$ (v2).  The
cross-sample bipartite transport finds substantially more structural signal
than within-sample spectral summaries.  This is the strongest separation
achieved by any spectral method in the symbiogenesis program.
\end{keyfinding}

\begin{keyfinding}
\textbf{All parent--child pairs rank above all pure unrelated pairs.}
Ranks 2, 3 (parent--child) $>$ ranks 8, 9, 10 (unrelated).  Ranks 6, 7
(parent--child) $>$ rank 8 (unrelated).  This qualitative ordering---related
pairs above unrelated pairs---is achieved for the first time by a spectral
method.
\end{keyfinding}

\subsection{Comparison to Prior Algorithms}

\begin{table}[H]
  \centering
  \tablestyle
  \caption{All symbiogenesis algorithms on CEPH~1463.}
  \label{tab:comparison}
  \begin{tabular}{lcccc}
    \toprule
    \rowcolor{table-header}
    \textcolor{white}{\textbf{Algorithm}} &
    \textcolor{white}{\textbf{Separation}} &
    \textcolor{white}{\textbf{Weakest Gap}} &
    \textcolor{white}{\textbf{Time}} &
    \textcolor{white}{\textbf{Status}} \\
    \midrule
    Rare-Run P90
      & $+$583\% & $+$300\% & 72.5\,s & \textbf{Gold} \\
    Module Stability v4
      & $+$4.90\% & $-$1.04\% & 405\,s & Fails \\
    \textbf{Spectral Ecology v3}
      & $\mathbf{+4.93\%}$ & $-$54.67\% & 742\,s & Fails \\
    Module Persistence v3
      & $+$3.36\% & $+$0.13\% & 1{,}085\,s & \textbf{PASSES} \\
    Coalition Transfer v3
      & $+$1.97\% & $+$0.06\% & 272\,s & \textbf{PASSES} \\
    Spectral Ecology v2
      & $+$1.55\% & $-$22.53\% & 379\,s & Fails \\
    \bottomrule
  \end{tabular}
\end{table}

\ymarginnote{Highest separation among all spectral methods; competitive with
Module Stability v4's $+4.90\%$.}

Spectral Ecology v3 achieves the highest separation of any symbiogenesis
algorithm except Rare-Run~P90, surpassing even Module Stability v4 ($+4.90\%$)
which held the previous record.  The weakest gap is large and negative
because the spouse pair score (1.477) is far above the lowest parent--child
score (0.930), but this reflects a specific, diagnosable failure mode rather
than a fundamental limitation.

\subsection{Component Analysis}

The three scoring components contribute different signals:

\begin{itemize}
  \item \textbf{Transport affinity ($T$)} has the widest range (1.11--2.64)
    and correlates with graph size similarity.  Father $\leftrightarrow$
    Mother scores highest ($T = 2.64$) because their graphs are nearly
    identical in size (444 vs.\ 443 modules), producing 1{,}400 mutual
    matches---nearly every module matches.
  \item \textbf{Topology preservation ($N$)} has a narrow range (0.33--0.41)
    and is the most size-invariant component.  Notably,
    Mat.GF~$\to$~Mother achieves the \emph{highest} $N$ score (0.411)
    despite ranking 7th overall---the matched modules have the most coherent
    neighborhood structure.
  \item \textbf{Unmatched penalty ($U$)} correlates with graph size overlap:
    pairs with similar $n$ have more matched modules and higher $U$.
\end{itemize}

\begin{keyfinding}
\textbf{Topology preservation is the most promising signal.}  It is the
least correlated with graph size and produces the highest score for a
parent--child pair (Mat.GF~$\to$~Mother, $N = 0.411$).  Future iterations
should increase its weight and refine the neighborhood consistency metric.
\end{keyfinding}


% ════════════════════════════════════════════════════════════════════
\section{Failure Analysis}
\label{sec:failure}

\subsection{The Indiscriminate Matching Problem}

The spouse pair (Father~$\leftrightarrow$~Mother) scores 1.477 with 1{,}400
mutual matches out of $\min(444, 443) = 443$ possible per side.  This means
each module in Father matches $\sim$3.2 modules in Mother and vice versa.
The matching is not selective---nearly everything matches.

\begin{limitation}
\textbf{When nearly every module finds a mate, transport affinity measures
graph size, not structural specificity.}  The score $T =
\sum K(i,j) / \sqrt{n_A n_B}$ grows with the number of matches.  For graphs
of similar size with similar feature distributions, mutual top-5 matching
produces $O(n)$ matches, making $T \propto \sqrt{n}$.
\end{limitation}

The fix is to make matching more selective:

\begin{enumerate}
  \item \textbf{Significance filtering.}  Only count a match if $K(i,j)$
    exceeds what would be expected by chance given the marginal feature
    distributions.
  \item \textbf{One-to-one constraint.}  Require each module to match at
    most one mate (Hungarian algorithm or Sinkhorn normalization), preventing
    the many-to-many explosion.
  \item \textbf{Topology-weighted affinity.}  Weight each match by its
    topology preservation score $N(i,j)$, so indiscriminate matches with
    incoherent neighborhoods contribute less.
\end{enumerate}

\subsection{The In-Law Problem}

Mat.GF~$\leftrightarrow$~Father (in-law, rank~4, score~1.011) outscores two
parent--child pairs (Pat.GF~$\to$~Father at 0.943, Mat.GF~$\to$~Mother
at 0.930).  This pair benefits from NA12891 (Mat.~GF) having the most modules
(449) and edges (1{,}410), producing many high-quality matches with
Father's large graph (444 modules).

The transport affinity ($T = 1.628$) is high because both graphs are large.
Topology preservation ($N = 0.372$) is moderate---not coherent enough to
indicate relatedness, but the large $T$ overwhelms it.

\subsection{What Would Fix It}

The diagnosis points to a single architectural change: \textbf{reweight the
composite to favor topology preservation over transport affinity}.

Current weights: $0.45T + 0.40N + 0.15U$.

If we use $0.20T + 0.60N + 0.20U$, the topology preservation signal---which
is the least size-correlated---would dominate.  Mat.GF~$\to$~Mother
($N = 0.411$) would rank higher, and the spouse pair's advantage from
indiscriminate matching ($T = 2.64$) would be attenuated.


% ════════════════════════════════════════════════════════════════════
\section{Discussion}
\label{sec:discussion}

\subsection{Cross-Sample Structure Is the Missing Ingredient}

The $3.2\times$ separation improvement ($+1.55\% \to +4.93\%$) validates the
central hypothesis: pair-specific cross-sample structure carries more
relatedness signal than within-sample spectral summaries.

\ymarginnote{The experiment confirmed the hypothesis: cross-sample
correspondence is the missing ingredient.}

This is consistent with the symbiogenesis thesis.  Inheritance does not just
produce similar ecosystem \emph{shapes}---it produces ecosystems with
\emph{specific shared substructures} that can be identified through soft
correspondence.  A parent and child share $\sim$50\% of haplotype blocks; the
modules derived from those blocks should be ecologically equivalent even if
their k-mer compositions differ.  The bipartite transport detects this
equivalence.

\subsection{The Role of Feature Engineering}

The 20-dimensional feature vector was designed to capture ecological character
without token identity.  The rank normalization removes absolute scale, so
a ``hub'' module (high degree, high HKS, core role) looks the same whether
it comes from a graph of 295 or 449 nodes.

The feature set is deliberately redundant: HKS and spectral position both
derive from the Laplacian eigenvectors, but HKS captures multi-scale diffusion
while spectral position captures cluster membership.  The clustering
coefficient adds local structure that eigenvectors miss.  The role one-hot
provides a discrete signal that is invariant to graph size.

Future work should ablate individual feature groups to determine which
carry the cross-sample signal and which add noise.

\subsection{Computation Profile}

\begin{table}[H]
  \centering
  \tablestyle
  \caption{Timing breakdown.}
  \label{tab:timing}
  \begin{tabular}{lr}
    \toprule
    \rowcolor{table-header}
    \textcolor{white}{\textbf{Phase}} &
    \textcolor{white}{\textbf{Time}} \\
    \midrule
    Module extraction (6 samples $\times$ 2 passes) & $\sim$600\,s \\
    Spectral decomposition (6 samples) & $\sim$140\,s \\
    Pairwise transport (10 pairs) & 1.5\,s \\
    \midrule
    \textbf{Total} & \textbf{741\,s} \\
    \bottomrule
  \end{tabular}
\end{table}

The module extraction is called twice per sample (once for the algorithm,
once inside \texttt{buildInteractionGraph})---this should be deduplicated to
halve prep time.  The pairwise transport phase is fast: 0.15\,s per pair for
$\sim$400$\times$400 distance matrices.

\subsection{Relationship to Optimal Transport}

The mutual top-$k$ matching is a simplified version of optimal transport.
Full optimal transport (Sinkhorn algorithm, Wasserstein distance) would find
the minimum-cost coupling between module distributions, but requires $O(n^2)$
memory and $O(n^2 \cdot \text{iterations})$ time.  For $n \sim 400$, this is
feasible---roughly $640$\,KB per pair and sub-second computation.

The advantage of full OT over mutual top-$k$: it naturally handles the
one-to-one constraint and produces a principled distance metric.  The
disadvantage: it loses the topology preservation component, which is the
most promising signal.

A hybrid---Sinkhorn-normalized coupling with topology-preservation
reweighting---may be the optimal next step.

\subsection{From v1 to v3: The Full Arc}

\begin{table}[H]
  \centering
  \tablestyle
  \caption{Evolution of Spectral Ecology across three versions.}
  \label{tab:arc}
  \begin{tabular}{llcc}
    \toprule
    \rowcolor{table-header}
    \textcolor{white}{\textbf{Version}} &
    \textcolor{white}{\textbf{Key Innovation}} &
    \textcolor{white}{\textbf{Separation}} &
    \textcolor{white}{\textbf{Insight}} \\
    \midrule
    v1 & Sequential adjacency
      & $-0.74\%$ & Graphs must be dense \\
    v2 & Co-occurrence edges
      & $+1.55\%$ & Within-sample has a ceiling \\
    \textbf{v3} & \textbf{Cross-sample transport}
      & $\mathbf{+4.93\%}$ & \textbf{Pair structure is the signal} \\
    \bottomrule
  \end{tabular}
\end{table}

Each version's failure directly motivated the next version's design.  v1's
graph collapse motivated v2's co-occurrence edges.  v2's within-sample ceiling
motivated v3's cross-sample transport.  v3's indiscriminate matching motivates
v4's selective, topology-weighted transport.


% ════════════════════════════════════════════════════════════════════
\section{Future Work}
\label{sec:future}

\subsection{Selective Transport}

The highest-priority improvement: make matching selective by requiring matches
to be significantly better than random.  Compute expected kernel values under
a null model where features are permuted, and only count matches that exceed
the null by a threshold.

\subsection{Topology-Dominant Scoring}

Increase topology preservation weight from 0.40 to 0.60 and decrease transport
affinity from 0.45 to 0.20.  The topology signal is the most size-invariant
component and produces the highest parent--child score for the hardest pair.

\subsection{One-to-One Matching}

Replace mutual top-$k$ with Sinkhorn-normalized optimal transport, which
naturally enforces one-to-one correspondence and produces a proper distance
metric.  This eliminates the indiscriminate many-to-many matching that
inflates the spouse score.

\subsection{Feature Ablation}

Systematically remove feature groups (intrinsic only, spectral only, HKS only,
role + HKS, full vector) to determine which carry the cross-sample signal.
The ablation will reveal whether the transport score is driven by genuine
ecological equivalence or by residual coverage correlation in the features.

\subsection{Coverage Normalization Ablation}

Compare raw support features vs.\ rank-normalized vs.\ support omitted
entirely.  If the score is dominated by support-correlated features, the
matching is still riding on sequencing depth.

\subsection{Persistent Homology Features}

Add topological descriptors (Betti numbers, persistence diagrams) to the node
feature vector.  Topological features are inherently size-invariant and could
provide coverage-orthogonal signal for cross-sample matching.


% ════════════════════════════════════════════════════════════════════
\section{Limitations}
\label{sec:limitations}

\begin{limitation}
\textbf{Negative weakest gap.}  The $-54.67\%$ gap is the worst in the
spectral family, driven by the spouse pair's indiscriminate matching
(1{,}400 mutual matches from 443--444 modules).
\end{limitation}

\begin{limitation}
\textbf{Transport affinity correlates with graph size.}  Pairs with similar
$n$ produce more matches and higher $T$.  The normalization by $\sqrt{n_A n_B}$
is insufficient.
\end{limitation}

\begin{limitation}
\textbf{Topology preservation has narrow range.}  $N$ spans 0.33--0.41,
providing limited discriminative power.  The neighborhood consistency metric
may need richer definition.
\end{limitation}

\begin{limitation}
\textbf{Double module extraction.}  Modules are extracted twice per sample
(once explicitly, once inside graph construction), doubling prep time.
\end{limitation}

\begin{limitation}
\textbf{No feature ablation.}  The 20-dimensional feature vector was designed
but not ablated.  Unknown which features carry signal vs.\ noise.
\end{limitation}

\begin{limitation}
\textbf{Single dataset.}  Six samples, 10 pairs, European ancestry.
\end{limitation}


% ════════════════════════════════════════════════════════════════════
\section{Conclusion}
\label{sec:conclusion}

Spectral Ecology v3 demonstrates that cross-sample structural correspondence
is the missing ingredient for nonlinear relatedness detection.  The shift
from within-sample spectral summaries (v2, $+1.55\%$) to pair-coupled
bipartite transport (v3, $+4.93\%$) produces a $3.2\times$ improvement in
mean separation---the largest gain from any single architectural change in the
symbiogenesis program.

\ymarginnote{The signal is in the coupling, not the summary.}

All four parent--child pairs rank above all four pure unrelated pairs.
The topology preservation component---measuring whether matched modules have
matched neighborhoods---produces the highest score for the hardest parent--child
pair and is the least correlated with graph size.

The spouse confound persists because the transport affinity component rewards
indiscriminate matching.  The fix is clear and the next iteration is
well-defined: selective, topology-weighted transport with one-to-one
correspondence constraints.

\begin{pullquote}
Related individuals share not just similar ecosystems, but ecosystems that
can be coherently aligned module by module, neighborhood by neighborhood.
The transport score measures this alignment.
\end{pullquote}

The symbiogenesis program has traversed three paradigms: token matching
(Module Persistence, Coalition Transfer), within-sample spectral invariants
(v2), and cross-sample structural correspondence (v3).  Each paradigm
exposed new signal and new failure modes.  The path forward is clearer than
ever: make the correspondence selective, weight it by topological coherence,
and the spectral ecosystem transport will become a genuine relatedness
detector.


% ════════════════════════════════════════════════════════════════════
\bibliographystyle{plain}
\bibliography{../references}

\end{document}
